\section{Primitive cycle sectors on general axial graphs}
\label{sec:primitive-cycles}

The unicyclic theorem has a coefficient-level extension to an arbitrary
finite axial graph.  This extension concerns one specified coupling
multidegree.  It is not an expansion of the full answer as a sum of
independent cycle functions.

Choose one reference orientation for every edge and formally complexify the
hopping parameters by treating $t_{xy}$ and $\bar t_{xy}$ as algebraically
independent for each reference-oriented edge $(x,y)$.  Define
$t_{yx}=\bar t_{xy}$ and $\bar t_{yx}=t_{xy}$.  The physical specialization
is $\bar t_{xy}=\overline{t_{xy}}$.  For an oriented simple cycle
$\eta=(x_1,\ldots,x_r,x_1)$, let
\begin{equation}
 W_\eta=\prod_{j=1}^r t_{x_jx_{j+1}},
 \qquad
 \bar W_\eta=\prod_{j=1}^r t_{x_{j+1}x_j},
 \qquad x_{r+1}=x_1.
 \label{eq:primitive-W-eta}
\end{equation}
The second monomial is the reverse orientation of the first.  Let
$\Pi_\eta$ project a polynomial onto the direct sum of the monomial spaces
spanned by
\begin{equation*}
 \{h_xW_\eta,h_x\bar W_\eta:x\in V\}.
\end{equation*}
We define
$\operatorname{Prim}_\eta m_{r+1}=\Pi_\eta m_{r+1}$.  Thus this sector
contains no additional hopping, longitudinal exchange, or field factor.
Finite dimensionality and the scaling $H\mapsto\lambda H$ imply that every
$m_q$ is a homogeneous polynomial of total parameter degree $q$.

\begin{proposition}[Primitive cycle coefficient]
\label{prop:primitive-cycle}
Let $G$ be a finite simple graph.  Equip it with the axial Hamiltonian
\eqref{eq:axial-H} and arbitrary irreducible site spins.  Let $A,B,C$ be
three distinct marked vertices and write
\begin{equation}
 \Mmod_\beta(AB,BC)=\sum_{q\geq0}\beta^q m_q.
 \label{eq:primitive-series}
\end{equation}
For every simple cycle $\eta$ of length $r$,
\begin{equation}
 \operatorname{Prim}_\eta m_{r+1}=0
 \label{eq:primitive-cycle-zero}
\end{equation}
unless $A,B,C\in V(\eta)$.  If the cycle contains the three marked
vertices, orient it so that they occur in the cyclic order
$A\to B\to C\to A$.  Then the formally complexified coefficient is
\begin{equation}
 \boxed{
 \operatorname{Prim}_\eta m_{r+1}
 =\frac{(-1)^rh_B}{\ii}
 \left(\prod_{x\in V(\eta)}\kappa_x\right)
 \left(W_\eta-\bar W_\eta\right).}
 \label{eq:primitive-cycle-formal}
\end{equation}
On the physical parameter slice,
\begin{equation}
 \boxed{
 \operatorname{Prim}_\eta m_{r+1}
 =2(-1)^r h_B
 \left(\prod_{x\in V(\eta)}\kappa_x\right)
 \operatorname{Im}W_\eta .}
 \label{eq:primitive-cycle-coefficient}
\end{equation}
No edge-disjointness assumption is imposed on the other cycles of $G$.
\end{proposition}

\begin{proof}
Fix $x\in V$ and extract the coefficients of $h_xW_\eta$ and
$h_x\bar W_\eta$.  Every parameter absent from these monomials is thereby set
to zero.  This includes all $J^z$ couplings, all fields except $h_x$, and all
hoppings outside $\eta$.  Retain both formal hopping orientations on $\eta$.
The surviving Hamiltonian is a tensor sum of the isolated cycle
Hamiltonian and, when $x$ lies outside the cycle, a one-site spectator
Hamiltonian.  Its Gibbs operator factorizes.  Normalized partial traces
remove every spectator factor exactly.

If $\eta$ contains $A,B,C$, the arbitrary-spin unicyclic theorem applies to
the physical slice of the isolated marked cycle.  Since that identity holds
for every complex hopping value, it extends polynomially to independent
$t$ and $\bar t$.  Projection onto $h_xW_\eta$ and $h_x\bar W_\eta$, for
each $x$ and both orientations, gives
Eq.~\eqref{eq:primitive-cycle-formal} and selects $h_B$.  A field on a
spectator component cannot enter the reduced state on $ABC$.

Suppose instead that $\eta$ omits at least one marked vertex.  If it omits
$C$, coefficient extraction gives
$\rho_{ABC}=\rho_{AB}\otimes\rho_C$.  The modular commutator reduces to
\begin{equation*}
 \ii\Tr\rho_{AB}
 [\log\rho_{AB},I_A\otimes\log\rho_B]=0.
\end{equation*}
The case in which $A$ is omitted is analogous.  If $B$ is omitted, the
lifted reduced logarithms are sums of one-site logarithms and commute.  These
three cases also cover cycles that omit two or all three marked vertices.
Polynomial coefficient extraction therefore proves
Eq.~\eqref{eq:primitive-cycle-zero} on the original graph.
\end{proof}

The proposition determines the complete first admissible coefficient when the
shortest hopping cycles are considered together.  Let $G_t$ be the graph
obtained by deleting the edges whose hopping amplitude is zero.

\begin{corollary}[Girth coefficient]
\label{cor:girth-onset}
Assume that $G_t$ contains a cycle and let $g$ be its girth.  Then
\begin{equation}
 m_q=0\quad(0\leq q\leq g)
 \label{eq:girth-lower-zero}
\end{equation}
and
\begin{equation}
 \boxed{
 m_{g+1}=2(-1)^g h_B
 \sum_{\substack{\eta\subseteq G_t\ \mathrm{simple}\,,\ |\eta|=g\\
                   A,B,C\in V(\eta)}}
 \left(\prod_{x\in V(\eta)}\kappa_x\right)
 \operatorname{Im}W_\eta .}
 \label{eq:girth-coefficient}
\end{equation}
Each unoriented cycle is counted once and is oriented in the marked cyclic
order.  Equation~\eqref{eq:girth-coefficient} remains valid when distinct
shortest cycles share edges or paths.
\end{corollary}

\begin{proof}
For any monomial in $m_q$, let $a_e$ and $b_e$ be the exponents of $t_e$ and
$\bar t_e$, and put $d_e=a_e-b_e$.  Local $U(1)$ invariance makes $d$ a
divergence-free integral flow.  Conjugation oddness excludes the zero flow.
Every nonzero integral circulation decomposes into oriented simple cycles.
Consequently,
\begin{equation*}
 \deg_{\mathrm{hop}}
 =\sum_e(a_e+b_e)
 \geq\sum_e|a_e-b_e|
 =\lVert d\rVert_1
 \geq g.
\end{equation*}
Equality requires one unit of circulation around one simple $g$-cycle.
Physical time-reversal oddness requires an odd field degree and therefore
adds at least one further factor.  This proves
Eq.~\eqref{eq:girth-lower-zero}.  At total degree $g+1$, no neutral return or
$J^z$ factor is possible.  The coefficient is consequently the sum of the
primitive sectors in Proposition~\ref{prop:primitive-cycle}, which
gives Eq.~\eqref{eq:girth-coefficient}.
\end{proof}

If $G_t$ is a forest, the local real-gauge criterion instead gives
$\Mmod_\beta(AB,BC)=0$ at every temperature.  Corollary
\ref{cor:girth-onset} is only a first-admissible-order statement.  It
does not establish a cactus formula and it does not imply cycle additivity
at later orders.

Two exact higher-order calculations delimit this point.  First, take five
spin-$1/2$ sites and two triangles $ABC$ and $BDE$ that share only $B$.
Set every $J^z$ to zero, set $h_B=1$ and all other fields to zero, and choose
unit-modulus hoppings with
\begin{equation}
 W_1=W_{ABC}=\ii,
 \qquad
 W_2=W_{BDE}=1.
\end{equation}
Exact coefficient extraction gives
\begin{equation}
 m_7=-\frac1{768}.
 \label{eq:cactus-counterexample}
\end{equation}
At hopping degree six, gauge flow restricts the mixed two-cycle sector on the
unit-modulus phase torus to a linear combination of
$\operatorname{Im}(W_1W_2)$ and
$\operatorname{Im}(W_1\overline{W_2})$.  The graph automorphism
$D\leftrightarrow E$ reverses the orientation of the second triangle while
fixing $A,B,C$, and therefore equates the two coefficients.  At $W_1=\ii$
and $W_2=1$, the single-cycle second harmonics
vanish.  Equation~\eqref{eq:cactus-counterexample} therefore fixes the mixed
sector as
\begin{equation}
 -\frac1{1536}
 \left[
 \operatorname{Im}(W_1W_2)
 +\operatorname{Im}(W_1\overline{W_2})
 \right]
 =-\frac1{768}\operatorname{Im}W_1\operatorname{Re}W_2.
 \label{eq:cactus-mixed-sector}
\end{equation}
The primary certificate also verifies the sign reversal under $W_2\mapsto
-W_2$.  The rational value in Eq.~\eqref{eq:cactus-counterexample} is
reproduced by a second exact implementation of the same formal-series method.
This is an implementation check, not an independent method.

Second, take the spin-$1/2$ triangle $A\!-\!B\!-\!D\!-\!A$ and attach the
marked site $C$ as a leaf at $B$.  Set $J^z=0$, $h_B=1$, and all other fields
to zero.  Choose
\begin{equation}
 t_{AB}=t_{BD}=t_{BC}=1,
 \qquad
 t_{DA}=\ii.
\end{equation}
The only cycle omits $C$.  The same primary exact certificate nevertheless
gives
\begin{equation}
 m_0=m_1=\cdots=m_7=0,
 \qquad
 m_8=-\frac7{368640}.
 \label{eq:off-target-cycle-counterexample}
\end{equation}
Neutral returns on the leaf make the off-target cycle visible at this later
order.  Equations~\eqref{eq:cactus-mixed-sector} and
\eqref{eq:off-target-cycle-counterexample} show that the natural
coefficient-level organization is by coupling multidegree, not by a sum of
independent functions of Wilson products.


\section{Why the axial hypothesis is necessary}
\label{sec:axial-boundary}

The onset theorems above require the absence of on-site transverse fields.
In this section we write $h_x^z\equiv h_x$ for the longitudinal field in
Eq.~\eqref{eq:axial-H}.  To locate the boundary of that hypothesis, add
\begin{equation}
 H_\perp=-\frac12\sum_x
 \left(g_xS_x^++\overline{g_x}S_x^-\right).
 \label{eq:transverse-field-H}
\end{equation}
With $S^+=S^x+\ii S^y$ and
$H_\perp=-\sum_x(h_x^xS_x^x+h_x^yS_x^y)$, our convention is
$g_x=h_x^x-\ii h_x^y$.  Under the local rotation
Eq.~\eqref{eq:gauge-unitary},
\begin{equation}
 t_{xy}\longmapsto
 e^{\ii(\varphi_x-\varphi_y)}t_{xy},
 \qquad
 g_x\longmapsto e^{\ii\varphi_x}g_x.
 \label{eq:transverse-gauge-action}
\end{equation}
The open-link phase $t_{AB}\overline{g_A}g_B$ is therefore gauge invariant.
It is not a Wilson product and it need not involve any closed path.
For a formally complexified monomial, let $a_{xy}$ and $b_{xy}$ be the
exponents of $t_{xy}$ and $\bar t_{xy}$ on each reference-oriented edge and
set $d_{xy}=a_{xy}-b_{xy}$.  Let $q_g(x)$ be the net power of $g_x$ minus
the net power of $\bar g_x$.  Use
\begin{equation*}
 (\operatorname{div}d)_x
 =\sum_{x\to y}d_{xy}-\sum_{y\to x}d_{yx}.
\end{equation*}
Gauge invariance now imposes
\begin{equation}
 \operatorname{div}d+q_g=0.
 \label{eq:transverse-flow-balance}
\end{equation}
Transverse insertions can therefore terminate an open hopping flow.  The
closed-circulation argument behind the girth bound no longer applies.

We first state the part of the order-seven coefficient that is independent
of the ambient graph.  The statement is a projected coefficient identity,
not a classification of the full $m_7$.

\begin{proposition}[Ambient-independent transverse-field sector]
\label{prop:local-transverse-sector}
Let $A,B,C$ be distinct vertices of a finite graph, with
$\{A,B\},\{B,C\}\in E$, and let these three vertices carry spin $1/2$.
Other vertices and couplings are arbitrary.  Formally regard
$t,\bar t,g,\bar g$ as independent variables and let
$\Pi_{\mathrm{loc}}$ project $m_7$ onto the two multidegrees
\begin{align}
 &J^z_{AB}h_B^z t_{BC}\bar t_{BC}t_{AB}\bar g_Ag_B,
 \nonumber\\
 &J^z_{AB}h_B^z t_{BC}\bar t_{BC}\bar t_{AB}g_A\bar g_B.
 \label{eq:local-transverse-multidegrees}
\end{align}
Then
\begin{equation}
 \Pi_{\mathrm{loc}}m_7
 =\frac{J^z_{AB}h_B^z t_{BC}\bar t_{BC}}{46080\ii}
 \left(t_{AB}\bar g_Ag_B-\bar t_{AB}g_A\bar g_B\right).
 \label{eq:local-transverse-formal}
\end{equation}
On the physical parameter slice this becomes
\begin{equation}
 \boxed{
 \Pi_{\mathrm{loc}}m_7
 =\frac{J^z_{AB}h_B^z|t_{BC}|^2}{23040}
 \operatorname{Im}\!\left(t_{AB}\overline{g_A}g_B\right).}
 \label{eq:local-transverse-physical}
\end{equation}
\end{proposition}

\begin{proof}
Every $m_q$ is a polynomial in the Hamiltonian parameters.  Extracting either
multidegree in Eq.~\eqref{eq:local-transverse-multidegrees} sets every absent
parameter to zero.  The resulting Hamiltonian is
$H_{ABC}\otimes I_{\mathrm{spectators}}$.  The Gibbs operator factorizes,
and the spectator dimension cancels in the partition function and in every
normalized partial trace.  The projected coefficient is therefore identical
to the coefficient on the isolated open path $A\!-\!B\!-\!C$.

The remaining three-qubit coefficient is evaluated by the self-contained
exact symbolic extraction described in Appendix~\ref{app:reproducibility},
over
\begin{equation*}
 \mathbb Q(\ii)[x_1,\ldots,x_7]
 \big/(x_1^2,\ldots,x_7^2).
\end{equation*}
It expands the normalized Gibbs state and both reduced logarithms through
order seven before taking the modular commutator.  The two coefficients are
\begin{align}
 \left[J^z_{AB}h_B^z t_{BC}\bar t_{BC}t_{AB}\bar g_Ag_B\right]m_7
 &=-\frac{\ii}{46080},
 \nonumber\\
 \left[J^z_{AB}h_B^z t_{BC}\bar t_{BC}\bar t_{AB}g_A\bar g_B\right]m_7
 &=\frac{\ii}{46080}.
 \label{eq:local-transverse-extracted-pair}
\end{align}
Discarding a square in the quotient ring is exact for these squarefree
targets.  Perturbative parameter degrees are nonnegative, so a term already
divisible by $x_j^2$ cannot later contribute to either target monomial.
Equation~\eqref{eq:local-transverse-formal} follows, and physical
specialization gives Eq.~\eqref{eq:local-transverse-physical}.
\end{proof}

The projection in Proposition~\ref{prop:local-transverse-sector} is the
ambient-independent statement.  The complete coefficient can contain other
local and nonlocal multidegrees.  A full formula is available in one sharply
delimited local model.

\begin{proposition}[Complete isolated-path coefficient]
\label{prop:complete-local-m7}
Consider exactly three qubits on the open path $A\!-\!B\!-\!C$.  Set
$h_A^z=h_C^z=0$ and retain only the eight parameter families
\begin{equation}
 t_{AB},\ t_{BC},\ g_A,\ g_B,\ g_C,\
 J^z_{AB},\ J^z_{BC},\ h_B^z.
 \label{eq:eight-local-parameters}
\end{equation}
Define the three gauge-invariant imaginary parts
\begin{align}
 I_{AB}&=\operatorname{Im}
 \left(t_{AB}\overline{g_A}g_B\right),
 \nonumber\\
 I_{BC}&=\operatorname{Im}
 \left(t_{BC}\overline{g_B}g_C\right),
 \nonumber\\
 I_{AC}&=\operatorname{Im}
 \left(t_{AB}t_{BC}\overline{g_A}g_C\right).
 \label{eq:local-phase-invariants}
\end{align}
Then
\begin{equation}
 m_0=m_1=\cdots=m_6=0
 \label{eq:local-lower-orders}
\end{equation}
and
\begin{align}
 m_7=\frac{h_B^z}{23040}\Bigg[{}
 &J^z_{AB}\left(|t_{BC}|^2-(J^z_{BC})^2\right)I_{AB}
 \nonumber\\
 {}+{}&J^z_{BC}\left(|t_{AB}|^2-(J^z_{AB})^2\right)I_{BC}
 \nonumber\\
 {}+{}&\frac{J^z_{AB}J^z_{BC}}2 I_{AC}
 \Bigg].
 \label{eq:complete-local-m7}
\end{align}
\end{proposition}

\begin{proof}
The computer-assisted exact proof uses the self-contained extraction
described in Appendix~\ref{app:reproducibility}.  Its exhaustive polynomial
expansion of
the Gibbs and reduced-logarithm series produces ten formal complex
monomials, paired into the five real sectors displayed in
Eq.~\eqref{eq:complete-local-m7}, and no other monomial through degree seven.
The separate exact-arithmetic engine checks the physical formula at five
specified nondegenerate Gaussian-rational points.  It tests an independent
implementation of the same formal-series method, not an independent method.
Exhaustive monomial extraction supplies the coefficient identity.
\end{proof}

If $J^z_{BC}=0$, Eq.~\eqref{eq:complete-local-m7}
reduces to Eq.~\eqref{eq:local-transverse-physical} and $g_C$ drops out.
Formula~\eqref{eq:complete-local-m7} is complete only under
Eq.~\eqref{eq:eight-local-parameters}.  Enabling $h_A^z$, $h_C^z$, or
additional ambient couplings permits further order-seven sectors.

Two finite rings show why this distinction matters for the axial onset law.
Take consecutive marked sites $A=0$, $B=1$, and $C=2$ on an oriented
spin-$1/2$ ring.  Choose
\begin{equation}
 t_{AB}=1+\ii,
 \qquad
 t_e=1\quad(e\ne AB),
 \qquad
 J^z_{AB}=h_B^z=g_A=g_B=1,
 \label{eq:transverse-ring-point}
\end{equation}
with every other $J^z$, longitudinal field, and transverse field equal to
zero.  On the six-cycle, the axial value obtained by setting $g_A=g_B=0$ is
\begin{equation}
 m_7^{\mathrm{axial}}=\frac1{2048}.
\end{equation}
At the point in Eq.~\eqref{eq:transverse-ring-point}, exact extraction from
the unprojected series through order seven instead gives
\begin{equation}
 m_7=\frac{49}{92160}
 =\frac1{2048}+\frac1{23040}.
 \label{eq:C6-transverse-contamination}
\end{equation}
The local open-link sector and the primitive six-cycle sector occur at the
same order.  The exact values here and on the seven-cycle below are certified
by the separate exact implementation summarized in
Appendix~\ref{app:reproducibility}.

On the seven-cycle at the same parameter point, the same exact extraction
gives
\begin{equation}
 m_0=m_1=\cdots=m_6=0,
 \qquad
 m_7=\frac1{23040}.
 \label{eq:C7-transverse-background}
\end{equation}
The primitive cycle multidegree is first admissible at order eight.  Thus the
full modular commutator is already nonzero before the winding contribution
can occur.  More generally, for a consecutive marked path on a spin-$1/2$
cycle of length $n\geq7$, the certified local multidegree in
Eq.~\eqref{eq:local-transverse-physical} has fixed order seven, while the
primitive cycle term has order $n+1$.  The separation between the two orders
therefore grows with $n$.

These examples do not show that the order-seven term is the first
transverse-field contribution in every model.  They do not classify the
complete transverse-field expansion, exclude lower-order sectors under other
parameter choices, or
give an order table for cycles of lengths three, four, and five.  The local
phase in Eq.~\eqref{eq:local-transverse-physical} is loop independent, but it
is not phase blind.  It measures the phase of $t_{AB}$ relative to the two
on-site transverse-field directions.

No arbitrary-spin closed formula for the transverse-field background is
asserted here.

One coefficient statement remains exact after transverse fields are added.
The primitive $h_BW_\gamma$ coefficient of the axial theorem is unchanged,
because extracting that multidegree sets every $g_x$ and $\bar g_x$ to zero.
It need not be the first term or the dominant term of the full modular
commutator.  Equations~\eqref{eq:C6-transverse-contamination} and
\eqref{eq:C7-transverse-background} prove that axiality is essential to the
leading-order theorem stated above.
